\chapter{Lynch Syndrome}

\section*{Definition and Overview}
Lynch syndrome, historically known as hereditary non-polyposis colorectal cancer (HNPCC), is an autosomal dominant genetic condition that significantly increases the risk of colorectal cancer and several extracolonic malignancies. These include cancers of the endometrium, ovaries, stomach, small bowel, urinary tract, brain, and skin. It is caused by germline mutations in DNA mismatch repair (MMR) genes or the EPCAM gene, leading to microsatellite instability (MSI). Identifying individuals with Lynch syndrome is critical for implementing intensive cancer screening, chemoprevention, and risk-reducing surgical measures.

\section*{Epidemiology}
Lynch syndrome is the most common cause of hereditary colorectal cancer, accounting for approximately 3\% of all colorectal cancer cases. In Australia and globally, it is estimated to affect around 1 in 270 to 1 in 300 individuals, though the vast majority of carriers remain undiagnosed. The average age of colorectal cancer diagnosis in individuals with Lynch syndrome is 45--50 years, significantly younger than in the general population. The risk of endometrial cancer is also high, particularly in female carriers, often exceeding their risk of colorectal cancer depending on the specific gene mutated.

\section*{Aetiology and Pathophysiology}
Lynch syndrome is caused by heterozygous germline mutations in one of four DNA mismatch repair (MMR) genes or epigenetic inactivation of MSH2 via EPCAM deletion.
\begin{itemize}
    \item \textbf{Mismatch repair (MMR) gene mutations:} mutations in MLH1, MSH2, MSH6, or PMS2 impair the cell's ability to correct single-base mismatches and insertion-deletion loops during DNA replication.
    \item \textbf{Second-hit inactivation:} tumorigenesis follows the somatic loss of the wild-type allele (the "second hit") in a target cell, leading to complete MMR deficiency.
    \item \textbf{Microsatellite instability (MSI):} accumulation of replication errors in repetitive DNA sequences (microsatellites), causing genomic instability and accelerated mutation rates in oncogenes and tumour suppressor genes.
    \item \textbf{Autosomal dominant inheritance:} offspring of a mutation carrier have a 50\% chance of inheriting the pathogenic variant.
    \item \textbf{EPCAM deletions:} deletion of the EPCAM gene leads to transcriptional inactivation of the neighbouring MSH2 gene in epithelial cells.
\end{itemize}

\section*{Clinical Features}
Lynch syndrome does not present with physical symptoms until a malignancy develops. However, a strong family history of cancer is the hallmark clinical feature.
\begin{itemize}
    \item Colorectal cancer presenting at an early age (often under 50 years) with symptoms like altered bowel habits, rectal bleeding, or unexplained iron-deficiency anaemia
    \item Endometrial cancer in women, often presenting with abnormal uterine bleeding (postmenopausal or intermenstrual bleeding)
    \item Synchronic or metachronic tumours (multiple primary cancers diagnosed at the same time or sequentially)
    \item A pedigree showing multiple family members affected by colorectal, endometrial, ovarian, or other Lynch-associated cancers
\end{itemize}

\section*{Differential Diagnosis}
It is essential to differentiate Lynch syndrome from other hereditary and sporadic colorectal cancer syndromes.
\begin{itemize}
    \item \textbf{Familial adenomatous polyposis (FAP):} characterized by hundreds to thousands of colorectal adenomatous polyps, caused by mutations in the APC gene.
    \item \textbf{MUTYH-associated polyposis (MAP):} an autosomal recessive polyposis syndrome causing increased colorectal cancer risk, mimicking mild FAP or Lynch syndrome.
    \item \textbf{Sporadic colorectal cancer with MSI:} about 15\% of sporadic colorectal cancers exhibit MSI, usually due to somatic promoter hypermethylation of the MLH1 gene, rather than a germline mutation.
    \item \textbf{Li-Fraumeni syndrome:} a rare hereditary cancer syndrome caused by TP53 mutations, presenting with sarcomas, breast cancer, and brain tumours at very young ages.
    \item \textbf{Cowden syndrome:} characterised by multiple hamartomas and increased risk of breast, thyroid, and endometrial cancers, caused by PTEN mutations.
\end{itemize}

\section*{When to Seek Medical Care}
Individuals with a personal or strong family history of early-onset colorectal cancer, endometrial cancer, or multiple Lynch-associated cancers should request a referral to a clinical genetics service or familial cancer clinic.

\textbf{Call 000 (or local emergency services) for:}
\begin{itemize}
    \item Severe, acute abdominal pain, vomiting, and inability to pass gas or stool (suggesting bowel obstruction)
    \item Massive, active rectal bleeding or collapse due to severe blood loss
    \item Sudden loss of consciousness or severe neurological deficits (suggesting a brain tumour or metastatic event)
    \item Severe, persistent shortness of breath or chest pain
\end{itemize}

\section*{Diagnosis and Investigations}
Diagnosis involves screening tumour tissue followed by germline genetic testing.
\begin{itemize}
    \item \textbf{Immunohistochemistry (IHC) on tumour tissue:} testing for the expression of the four MMR proteins (MLH1, MSH2, MSH6, PMS2) to detect loss of one or more proteins.
    \item \textbf{Microsatellite instability (MSI) testing:} PCR-based analysis of tumour DNA to detect characteristic changes in microsatellite lengths.
    \item \textbf{BRAF V600E mutation and MLH1 promoter methylation analysis:} performed on tumours lacking MLH1 expression; presence of BRAF mutation or promoter methylation strongly suggests a sporadic, non-hereditary origin.
    \item \textbf{Germline genetic testing:} DNA sequencing of MLH1, MSH2, MSH6, PMS2, and EPCAM from a blood or saliva sample to identify a pathogenic germline variant, providing definitive diagnosis.
\end{itemize}

\section*{Management}
Management focuses on intensive surveillance, risk-reducing surgery, and chemoprevention.
\begin{itemize}
    \item \textbf{Regular colonoscopy:} performed every 1 to 2 years starting from age 20 to 25 (or earlier depending on the gene), with removal of any precancerous adenomas.
    \item \textbf{Hysterectomy and bilateral salpingo-oophorectomy:} offered as risk-reducing surgery to women who have completed childbearing to prevent endometrial and ovarian cancers.
    \item \textbf{Aspirin chemoprevention:} daily aspirin has been shown in clinical trials to significantly reduce the risk of colorectal cancer in Lynch syndrome carriers.
    \item \textbf{Surveillance for other cancers:} upper endoscopy every 2 to 3 years to screen for gastric and small bowel cancers, and urinalysis for urothelial cancers in selected high-risk families.
    \item \textbf{Surgical resection of malignancies:} if colorectal cancer develops, subtotal or total colectomy is often preferred over segmental resection due to the high risk of metachronic tumours.
\end{itemize}

\section*{Complications}
\begin{itemize}
    \item Development of multiple primary cancers (metachronic colorectal or extracolonic tumours)
    \item Adverse effects of intensive surveillance, such as bowel perforation during colonoscopy
    \item Psychological distress, anxiety, and depression related to cancer risk and genetic transmission to children
    \item Surgical complications from risk-reducing hysterectomy or major bowel resection
    \item Financial and insurance challenges associated with a documented genetic predisposition
\end{itemize}

\section*{Prognosis and Outcomes}
With active surveillance and preventative management, the prognosis for individuals with Lynch syndrome is excellent. Colonoscopic surveillance reduces colorectal cancer mortality by over 60\%. Although carriers have a high lifetime risk of cancer, early detection allows for curative treatment. The prognosis of Lynch-associated colorectal cancers is generally better than that of sporadic cancers, partly due to the prominent immune response triggered by mismatch-repair deficient tumour cells.

\section*{Prevention and Self-Care}
\begin{itemize}
    \item Adhere strictly to the recommended colonoscopy and screening schedule.
    \item Discuss the option of daily preventive aspirin with a specialist or GP.
    \item Maintain a healthy weight and engage in regular physical activity to reduce overall cancer risk.
    \item Avoid smoking and limit alcohol consumption, as these can exacerbate colorectal cancer risk.
    \item Communicate genetic test results to first-degree relatives so they can undergo predictive testing and early screening.
\end{itemize}

\section*{Sources and Further Reading}
The following reputable organisations provide further information on Lynch syndrome, genetic testing, and cancer surveillance.
\begin{itemize}
    \item Cancer Council Australia. \textit{Lynch syndrome screening and management guidelines.} https://www.cancer.org.au/
    \item eviQ (Cancer Institute NSW). \textit{Lynch syndrome genetic testing and surveillance protocols.} https://www.eviq.org.au/
    \item Mayo Clinic. \textit{Lynch syndrome symptoms, genetics, and prevention.} https://www.mayoclinic.org/diseases-conditions/lynch-syndrome/
    \item Genetics Home Reference (MedlinePlus). \textit{Lynch syndrome genetic details.} https://medlineplus.gov/genetics/condition/lynch-syndrome/
    \item NICE. \textit{Testing strategies for Lynch syndrome in people with colorectal cancer.} https://www.nice.org.uk/guidance/dg27
\end{itemize}
